\documentclass[12pt, a4paper]{article}

% ============================================================
%  SwiftLaTeX 示例项目
%  使用 XeLaTeX 编译，支持中英文混排
% ============================================================

% 中文支持
\usepackage{ctex}

% 页面布局
\usepackage[margin=2.5cm]{geometry}

% 数学公式
\usepackage{amsmath, amssymb}

% 图片
\usepackage{graphicx}

% 超链接
\usepackage{hyperref}
\hypersetup{
    colorlinks=true,
    linkcolor=blue,
    citecolor=blue,
    urlcolor=blue
}

% 参考文献
\usepackage[numbers]{natbib}

% 代码高亮
\usepackage{listings}
\lstset{
    basicstyle=\ttfamily\small,
    frame=single,
    breaklines=true
}

% ============================================================

\title{\textbf{SwiftLaTeX 示例文档} \\ \large 一个简单的 \LaTeX{} 入门项目}
\author{SwiftLaTeX 用户}
\date{\today}

\begin{document}

\maketitle

\begin{abstract}
这是一个 SwiftLaTeX 的示例项目，用于演示编辑器的基本功能。
本文档包含中英文混排、数学公式、代码块、表格和参考文献等常见 \LaTeX{} 元素，
帮助你快速上手 SwiftLaTeX。
\end{abstract}

\tableofcontents
\newpage

% ============================================================
\section{欢迎使用 SwiftLaTeX}
% ============================================================

SwiftLaTeX 是一个基于浏览器的 \LaTeX{} IDE，提供以下核心功能：

\begin{itemize}
    \item \textbf{代码编辑器} — 基于 CodeMirror 6，支持语法高亮和自动补全
    \item \textbf{实时编译预览} — 修改后自动编译，实时查看 PDF 输出
    \item \textbf{AI 助手} — 支持多种 AI 模型，帮你修复错误、润色文章
    \item \textbf{评论系统} — 在具体行添加评论，方便协作讨论
    \item \textbf{TODO 管理} — 跟踪任务进度
\end{itemize}

你可以通过快捷键 \texttt{Ctrl+S}（或 \texttt{Cmd+S}）保存文件，
编辑器会在 2 秒无操作后自动保存。

% ============================================================
\section{数学公式}
% ============================================================

\LaTeX{} 的数学排版能力是它最强大的特性之一。

\subsection{行内公式}

著名的欧拉公式 $e^{i\pi} + 1 = 0$ 被誉为数学中最美的公式。
质能方程 $E = mc^2$ 揭示了质量与能量的等价关系。

\subsection{行间公式}

高斯分布（正态分布）的概率密度函数：
\begin{equation}
    f(x) = \frac{1}{\sigma\sqrt{2\pi}} \exp\left(-\frac{(x-\mu)^2}{2\sigma^2}\right)
    \label{eq:gaussian}
\end{equation}

贝叶斯定理：
\begin{equation}
    P(A \mid B) = \frac{P(B \mid A) \cdot P(A)}{P(B)}
    \label{eq:bayes}
\end{equation}

矩阵运算示例：
\begin{equation}
    \mathbf{A} = \begin{pmatrix}
        a_{11} & a_{12} & a_{13} \\
        a_{21} & a_{22} & a_{23} \\
        a_{31} & a_{32} & a_{33}
    \end{pmatrix}
\end{equation}

% ============================================================
\section{代码展示}
% ============================================================

以下是一个简单的 Python Hello World 程序：

\begin{lstlisting}[language=Python, caption={Python 示例}]
def hello(name: str) -> str:
    """Say hello to someone."""
    return f"Hello, {name}! Welcome to SwiftLaTeX."

if __name__ == "__main__":
    print(hello("World"))
\end{lstlisting}

% ============================================================
\section{表格}
% ============================================================

表~\ref{tab:comparison} 对比了几种常见的 \LaTeX{} 编辑器。

\begin{table}[htbp]
    \centering
    \caption{LaTeX 编辑器对比}
    \label{tab:comparison}
    \begin{tabular}{lccc}
        \hline
        \textbf{编辑器} & \textbf{本地运行} & \textbf{AI 助手} & \textbf{免费} \\
        \hline
        SwiftLaTeX  & \checkmark & \checkmark & \checkmark \\
        Overleaf    &            & \checkmark &            \\
        TeXstudio   & \checkmark &            & \checkmark \\
        VS Code     & \checkmark &            & \checkmark \\
        \hline
    \end{tabular}
\end{table}

% ============================================================
\section{交叉引用}
% ============================================================

如公式~\eqref{eq:gaussian} 所示，高斯分布是统计学中最重要的分布之一。
公式~\eqref{eq:bayes} 是贝叶斯推断的基础。

关于深度学习的更多内容，可以参考 Goodfellow 等人的经典教材~\cite{goodfellow2016deep}。
Transformer 架构~\cite{vaswani2017attention} 彻底改变了自然语言处理领域。

% ============================================================
\section{总结}
% ============================================================

恭喜你成功编译了第一个 SwiftLaTeX 项目！接下来你可以：

\begin{enumerate}
    \item 在 \texttt{projects/} 目录下创建自己的项目
    \item 尝试使用 AI 助手帮你润色文章
    \item 给特定行添加评论
    \item 在 TODO 面板中管理任务
\end{enumerate}

如果遇到问题，欢迎在 GitHub 上提 Issue：\\
\url{https://github.com/ChangXiang-SCU/SwiftLaTeX}

% ============================================================
% 参考文献
% ============================================================

\bibliographystyle{plainnat}
\bibliography{references}

\end{document}
